Sources : études citées, présentées dans le complément « Les élasticités estimées sur données françaises » ; chaque chiffre, avec sa page et son tableau, figure dans donnees/elasticites\_france.csv.\par Lecture : chez Lefebvre et al., les deux élasticités sont estimées ensemble sur les mêmes 2,85 millions de foyers, dont le revenu fiscal de référence dépasse 30 000 \euro{} et qui déclarent chaque année des revenus des deux sortes ; les revenus d'activité comprennent les pensions, qui ne réagissent pas. La barre claire couvre les quinze spécifications publiées : dans chacune, l'élasticité des revenus du capital dépasse celle des revenus d'activité d'au moins 0,34 ; les intervalles de confiance sont plus étroits que les points. Chez Bach et al., l'élasticité vient de la comparaison, sur plusieurs années, des assujettis à l'ISF touchés par la réforme et des autres ; aucun écart-type n'est publié. En 2013, la détention des sociétés est mal mesurée, ce qui, selon les auteurs, surestime la réaction des autres foyers et sous-estime celle des dirigeants.
